a hyperfield. Specifically, by \(-xyz\) we mean the additive inverse of the product \(xyz\), and by \(-xyz\boxplus -xyz\boxplus -xyz\) we mean three separate summands equal to \(-xyz\).

We prove the following theorem.

\begin{theorem}\label{thm:main}
Let \(\mathbb H\) be a hyperfield and let \(x,y,z\in\mathbb H\). If
\[
0\in x\boxplus y\boxplus z,
\]
then
\[
0\in x^{3}\boxplus y^{3}\boxplus z^{3}\boxplus (-xyz)\boxplus(-xyz)\boxplus(-xyz).
\]
In particular, the answer to the question is \emph{yes}.
\end{theorem}

\section{Hyperfield preliminaries}

We use the standard axioms of a Krasner hyperfield. Thus \((\mathbb H,\boxplus,0)\) is a canonical commutative hypergroup, \((\mathbb H\setminus\{0\},\cdot,1)\) is an abelian group, and multiplication distributes over the multivalued addition. In particular:

\begin{itemize}
\item every element \(a\in\mathbb H\) has a unique additive inverse \(-a\), and
\[
0\in a\boxplus(-a);
\]
\item for \(a,b\in\mathbb H\),
\[
0\in a\boxplus b \quad\Longleftrightarrow\quad b=-a;
\]
\item for all \(a,b,c\in\mathbb H\),
\[
a(b\boxplus c)=ab\boxplus ac,
\]
where multiplication of a subset is understood elementwise.
\end{itemize}

For a finite sequence \(a_1,\dots,a_n\in\mathbb H\), the iterated hypersum is defined recursively by
\[
a_1\boxplus\cdots\boxplus a_n
=
\bigcup_{u\in a_2\boxplus\cdots\boxplus a_n} a_1\boxplus u.
\]
By associativity and commutativity, the order of the summands and the placement of parentheses do not affect the resulting subset of \(\mathbb H\).

We shall also use the following elementary observation.

\begin{lemma}\label{lem:witness}
Let \(\mathbb H\) be a hyperfield and \(x,y,z\in\mathbb H\). Then
\[
0\in x\boxplus y\boxplus z
\]
if and only if
\[
-z\in x\boxplus y.
\]
\end{lemma}

\begin{proof}
By associativity,
\[
0\in x\boxplus y\boxplus z
\]
means that there exists \(u\in x\boxplus y\) such that
\[
0\in u\boxplus z.
\]
By the inverse axiom for hyperfields, this is equivalent to \(u=-z\). Hence
\[
-z\in x\boxplus y.
\]
The converse is immediate, since if \(-z\in x\boxplus y\), then
\[
0\in (-z)\boxplus z\subseteq x\boxplus y\boxplus z.
\]
\end{proof}

We also recall the standard interpretation of integer coefficients in a hyperfield. If a monomial \(m\) occurs with coefficient \(+k\), we interpret it as \(k\) copies of the element obtained by evaluating \(m\). If it occurs with coefficient \(-k\), we interpret it as \(k\) copies of the additive inverse of the element obtained by evaluating \(m\). Thus, for example, the polynomial
\[
X^{3}+Y^{3}+Z^{3}-3XYZ
\]
evaluates at \((x,y,z)\) to
\[
x^{3}\boxplus y^{3}\boxplus z^{3}
\boxplus(-xyz)\boxplus(-xyz)\boxplus(-xyz).
\]
This is exactly the expression appearing in the question.

\section{A standard divisibility theorem for roots over hyperfields}

We shall use the following standard theorem, due in this formulation to Viro. It says that the usual divisibility relation between polynomials gives the expected implication for roots over arbitrary hyperfields.

\begin{theorem}[Viro, linear-factor divisibility theorem]\label{thm:viro}
Let \(L\in\mathbb Z[X_1,\dots,X_n]\) be a homogeneous linear form whose nonzero coefficients are \(\pm1\). Let \(P\in\mathbb Z[X_1,\dots,X_n]\) be homogeneous, and suppose that
\[
P=LQ
\]
in \(\mathbb Z[X_1,\dots,X_n]\) for some \(Q\in\mathbb Z[X_1,\dots,X_n]\). Let \(\mathbb H\) be a hyperfield and let \(a=(a_1,\dots,a_n)\in\mathbb H^n\). If \(0\in L(a)\), then
\[
0\in P(a).
\]
\end{theorem}

A proof of this theorem, in slightly different notation, is given in Viro's work on hyperfields and dequantization; see \cite{Viro}, Section~4, where roots of polynomials over hyperfields are defined exactly as above and the basic divisibility property is proved. The result is often summarized by saying that if a polynomial is a multiple of a linear form in the usual polynomial ring, then every hyperfield zero of the linear form is a hyperfield zero of the multiple.

For convenience, we spell out the interpretation of \(L(a)\) in a simple case. If
\[
L=X+Y-T,
\]
then
\[
L(x,y,t)=x\boxplus y\boxplus(-t).
\]
Thus \(0\in L(x,y,t)\) means
\[
0\in x\boxplus y\boxplus(-t),
\]
or equivalently, by Lemma~\ref{lem:witness},
\[
t\in x\boxplus y.
\]

\section{Proof of the main theorem}

Assume that
\[
0\in x\boxplus y\boxplus z.
\]
By Lemma~\ref{lem:witness}, this is equivalent to
\[
-z\in x\boxplus y.
\]
Set
\[
t=-z.
\]
Then
\[
t\in x\boxplus y.
\]
Equivalently,
\[
0\in x\boxplus y\boxplus(-t).
\]

Consider the homogeneous linear form
\[
L=X+Y-T.
\]
Then the point \((x,y,t)\in\mathbb H^3\) is a zero of \(L\), because
\[
0\in x\boxplus y\boxplus(-t).
\]

Now consider the polynomial
\[
P=X^{3}+Y^{3}-T^{3}+3XYT.
\]
We claim that
\[
P=LQ
\]
in \(\mathbb Z[X,Y,T]\), where
\[
Q=T^{2}+XT+YT+X^{2}-XY+Y^{2}.
\]
Indeed,
\[
\begin{aligned}
&(X+Y-T)\bigl(T^{2}+XT+YT+X^{2}-XY+Y^{2}\bigr)\\
&=X^{3}+Y^{3}-T^{3}+3XYT.
\end{aligned}
\]
This identity is verified by direct expansion.

By Theorem~\ref{thm:viro}, since \(0\in L(x,y,t)\), we obtain
\[
0\in P(x,y,t).
\]
Using our convention for integer coefficients, this means
\[
0\in x^{3}\boxplus y^{3}\boxplus(-t^{3})
\boxplus xyt\boxplus xyt\boxplus xyt.
\]

It remains to rewrite this in terms of \(x,y,z\). Since \(t=-z\), we have
\[
t^{3}=(-z)^{3}=-z^{3}.
\]
Therefore
\[
-t^{3}=z^{3}.
\]
Also,
\[
xyt=xy(-z)=-xyz.
\]
Hence
\[
x^{3}\boxplus y^{3}\boxplus(-t^{3})
\boxplus xyt\boxplus xyt\boxplus xyt
=
x^{3}\boxplus y^{3}\boxplus z^{3}
\boxplus(-xyz)\boxplus(-xyz)\boxplus(-xyz).
\]
Thus
\[
0\in x^{3}\boxplus y^{3}\boxplus z^{3}
\boxplus(-xyz)\boxplus(-xyz)\boxplus(-xyz).
\]

This completes the proof of Theorem~\ref{thm:main}. \qed

\section{Remark}

The use of Theorem~\ref{thm:viro} is essential. Over a hyperfield, addition is multivalued, so one cannot simply argue that the ordinary polynomial identity
\[
X^{3}+Y^{3}+Z^{3}-3XYZ
=
(X+Y+Z)(X^{2}+Y^{2}+Z^{2}-XY-XZ-YZ)
\]
implies the desired hyperfield statement by a purely formal manipulation. The theorem of Viro supplies precisely the needed invariance under the polynomial cancellations occurring in the displayed factorization.

\begin{thebibliography}{9}

\bibitem{Viro}
O.~Viro,
\emph{Hyperfields for tropical algebra I. Hyperfields and dequantization},
preprint, arXiv:1006.0807, 2010.

\end{thebibliography}

\end{document}